% Self-contained mathematical subsection; include after theorem/macros setup.
\section{The dodecahedron: a finite exact certificate}
Write $p=(1+\sqrt5)/2$, so $p^2=p+1$.  The coordinate model is the convex
hull of
\[
 (\pm1,\pm1,\pm1),\qquad
 (0,\pm(p-1),\pm p),\qquad
 (\pm(p-1),\pm p,0),\qquad
 (\pm p,0,\pm(p-1)).
\]
Its edge length is $2/p$.  Its facets have outward normals
$(0,\pm p,\pm1),(\pm1,0,\pm p),(\pm p,\pm1,0)$ and support constant $p^2$.
Every facet is a regular pentagon.  These assertions are elementary finite
coordinate checks; the supplied verifier checks all facet incidences and all
30 edges exactly.

\subsection{The attaining section}
The plane $p y+z=0$ cuts out a regular decagon of side length $1$ and
circumradius $p$.  Five consecutive vertices are
\[
 (p,0,0),\quad
 \left(\frac{p^2}{2},\frac12,-\frac p2\right),\quad
 \left(\frac12,\frac p2,-\frac{p^2}{2}\right),\quad
 \left(-\frac12,\frac p2,-\frac{p^2}{2}\right),\quad
 \left(-\frac{p^2}{2},\frac12,-\frac p2\right),
\]
followed by their negatives in the same order.  Each lies on a dodecahedral
edge midpoint.  Consecutive vertices have squared distance $1$ and all
vertices have squared distance $p^2$ from the origin.  The listed cyclic
order and convexity therefore give a regular decagon without requiring
trigonometric functions.  Its apothem is
$\sqrt{p^2-1/4}=\frac12\sqrt{4p+3}$, so triangulation from the centre gives
\[
 M=\frac52\sqrt{4p+3}=\frac52\sqrt{5+2\sqrt5},
 \qquad M^2=\frac{25}{4}(4p+3).
\]
For a dodecahedron of edge length $a$, scaling by $ap/2$ gives
\[
 \boxed{\ A_{\max}=\frac{5}{16}\sqrt{130+58\sqrt5}\,a^2\ }.
\]

\subsection{Reduction to three cones of normals}
By the central-section lemma, it suffices to consider planes through the
origin.  The remaining reduction uses only a finite reflection group.
Set
\[
 F=(0,p,1),\qquad E=(1,p,p+1),\qquad V=(1,1,1).
\]
These rays span a reflection chamber.  Inward normals to its three walls are
\[
 \alpha_0=(-1,p,1-p),\quad
 \alpha_1=(1-p,-1,p),\quad
 \alpha_2=(p+1,1,-p).
\]
Here $\alpha_i$ is orthogonal to the two other rays and has positive dot
product with ray $i$.  Reflection in each wall permutes the twenty vertices.
Closing these three reflections under composition produces precisely
120 orthogonal maps, as checked by exact arithmetic.

Here is a constructive proof that their chambers cover all normals.  Put
$c=F+E+V$.  From the finite list choose $g$ maximizing $n\cdot(gc)$, and put
$n'=g^{-1}n$.  For every wall reflection $s_i$, maximality implies
$n'\cdot c\ge n'\cdot(s_i c)$.  Since
\[
 c-s_i c=\frac{2(c\cdot\alpha_i)}{\alpha_i\cdot\alpha_i}\alpha_i
 \quad\hbox{and}\quad c\cdot\alpha_i>0,
\]
we obtain $n'\cdot\alpha_i\ge0$.  Thus $n'$ is in the positive cone generated
by $F,E,V$.  This argument is a bounded maximum over a finite, explicitly
verified list; it does not use a continuous optimization theorem.

Write $n'=aF+bE+cV$ with $a,b,c\ge0$.  Only two vertex-sign planes cross
this cone.  They are
\[
 n'\cdot(-1,1,-1)=(p-1)a-2b-c=0,
 \qquad
 n'\cdot(-p,0,p-1)=(p-1)a-c=0.
\]
Put
\[
 T=2F+(p-1)E=(p-1,2p+1,p+2),\qquad
 S=F+(p-1)V=(p-1,2p-1,p).
\]
The chamber splits into the three simplicial cones
\[
 C_0=\operatorname{cone}(F,T,S),\quad
 C_1=\operatorname{cone}(T,E,S),\quad
 C_2=\operatorname{cone}(E,V,S).
\]
For completeness, put $d=(p-1)a$.  If $2b+c\le d$, the coordinates in
$C_0$ are $(a-p(2b+c),pb,pc)$.  If $c\le d\le2b+c$, the coordinates in
$C_1$ are $((a-pc)/2,(2b+c-d)/2,pc)$.  If $d\le c$, the coordinates in
$C_2$ are $(b,c-d,a)$.  Every displayed coordinate is nonnegative in its
stated case, and direct substitution reconstructs $n'$.  For arithmetic
inputs these case distinctions are decidable; the extension to represented
real planes is by the common quantitative approximation lemma.

\subsection{Rational area formulas}
Let $(x,y,z)$ denote the normal $n'$; throughout these cones $z>0$.
Define
\[
 q_1=(2-p)y+(p-1)z,\qquad q_2=py-(p-1)z,
\]
\[
 A=(p-1)(2yz-z^2)-x^2,
\]
\[
 B=(5-3p)y^2+(12p-18)yz+(14-9p)z^2-x^2,
\]
\begin{align*}
 Q={}&x^3+(2-p)x^2y+5(p-1)x^2z+(2p-6)xyz\\
    &+(7-3p)xz^2+2(1-p)y^2z-3yz^2+pz^3.
\end{align*}
For $n'\in C_i$, the section area satisfies
\[
 \operatorname{area}(P\cap n'^\perp)^2
       =(x^2+y^2+z^2)\frac{N_i(x,y,z)^2}{D_i(x,y,z)^2},
\]
where
\begin{align*}
 N_0&=2p^3AB,&
 D_0&=z(q_1^2-x^2)(q_2^2-x^2),\\
 N_1&=-2p^3Q,&
 D_1&=z\bigl(x+py+(p+1)z\bigr)(x+q_1)(q_2-x),\\
 N_2&=4p^2\bigl(x^2+y^2+z^2+2p(xy+xz+yz)\bigr),&
 D_2&=\bigl(x+py+(p-1)z\bigr)\bigl(x+py+(p+1)z\bigr)(x+q_1).
\end{align*}
The middle formula is used away from ray $E$; the last formula applies on
that ray.  No other denominator singularities occur on the stated domains.

Here is a derivation that also specifies a transparent independent way of
checking the formulas.  Number vertices in lexicographic sign order, first
the eight cube vertices, then the four vertices of each of the other three
families in the order displayed above.  If an edge has endpoints $v_i,v_j$,
its intersection with the plane is
\[
 w_{ij}=\frac{(n'\cdot v_i)v_j-(n'\cdot v_j)v_i}
                   {n'\cdot(v_i-v_j)}.
\]
The cyclic edge lists for the three cones are respectively
\begin{align*}
 C_0:;&(5,9),(5,19),(18,19),(6,18),(6,10),\\
      &(2,10),(2,16),(16,17),(1,17),(1,9);\\
 C_1:;&(5,14),(18,19),(6,18),(6,10),\\
      &(2,13),(16,17),(1,17),(1,9);\\
 C_2:;&(5,14),(4,18),(6,10),(2,13),(3,17),(1,9).
\end{align*}
The vertex signs at the cone rays verify that these are exactly the cut
edges.  Consecutive cut edges belong to a common pentagonal facet, giving
the indicated cyclic order.  The projected shoelace formula gives
\[
 \frac{\operatorname{area}}{\sqrt{x^2+y^2+z^2}}
 =\left|\frac1{2z}\sum_k
   \bigl((w_k)_x(w_{k+1})_y-(w_k)_y(w_{k+1})_x\bigr)\right|.
\]
Substitute the edge-intersection formula, clear its linear denominators,
and use $p^2=p+1$.  This yields exactly $|N_i/D_i|$ above.  The supplied
verifier checks the complete cleared polynomial identity, rather than
sampling orientations or trusting a symbolic simplification.  Boundary
vertex intersections merely duplicate listed polygon vertices and give
zero-length sides.  The only case in which a listed edge is entirely in
the section plane is the explicitly excluded ray $E$ in $C_1$, already
handled by $C_2$.

\subsection{The global upper bound by positive coefficients}
For each $i$, substitute
\[
 (x,y,z)=uR_{i0}+vR_{i1}+wR_{i2},\qquad u,v,w\ge0,
\]
using the ordered rays of $C_i$.  Form the homogeneous polynomial
\[
 H_i(u,v,w)=25(4p+3)D_i^2
       -4(x^2+y^2+z^2)N_i^2.
\]
Every nonzero coefficient of $H_i$ is strictly positive in $\mathbb Q(p)$.
Here is an exact small certificate summary:
\[
\begin{array}{c|c|c|c}
 i&\deg H_i&\text{nonzero coefficients}&\text{smallest certified lower bound}\\\hline
 0&10&63&28848/25\\
 1&8&42&50163872/625\\
 2&6&28&37139298/625
\end{array}
\]
The coefficient test requires only the rational enclosure
\[
 L=\frac{161803}{100000}<p<\frac{161804}{100000}=U,
\]
proved by $L^2-L-1<0<U^2-U-1$ and $L>1$.  For a coefficient $a+bp$, use
$a+bL$ when $b\ge0$, and $a+bU$ when $b<0$.  The last table gives the
minimum of these rational lower bounds, so it certifies all signs without
numerical approximations.  Every coefficient pair $(a,b)$ is included in
\texttt{dodecahedron\_certificate.json}; the optional coefficient appendix
prints the same complete finite data.  The verifier recomputes the
polynomials from $N_i,D_i$ and checks them coefficient by coefficient.
For example the first terms of $H_0$ are
\[
 (-38400+25600p)u^8v^2+
 (-38400+25600p)u^8vw+
 (-102400+64000p)u^8w^2.
\]
All omitted terms in this illustration, but none in the supplied certificate,
have the same verified positivity property.

Since $u,v,w\ge0$, we have $H_i\ge0$.  Each $D_i$, after the same substitution,
has nonnegative coefficients and strictly positive pure-power coefficients;
for $D_1$ the pure $v^d$ coefficient vanishes, exactly at its excluded ray.
Thus $D_i>0$ on its domain.  Dividing the certified inequality by $4D_i^2$
yields area squared at most $25(4p+3)/4=M^2$.  Nonnegativity of area now gives
area at most $M$, and the exhibited decagon attains it.  This is a global
bound, not an inference from the highly symmetric candidate alone.

\subsection{What is primitive recursive here}
The verification is a fixed finite calculation over pairs of rational
numbers.  Rational arithmetic, polynomial coefficient arrays, convolution,
finite permutations, squared-distance comparisons, and finite maxima have
primitive recursive implementations.  A proof of the coefficient checker
uses only quantifier-free induction on bounded lists and arithmetic
identities.  Each accepted positivity certificate is a displayed sum of
products of nonnegative quantities; its soundness needs no compactness,
excluded middle for real numbers, or unbounded search.  Interpreting the
geometry with arithmetic coordinates gives a direct iPRA proof of the
sharp squared-area inequality.  General represented inputs are handled by
the quantitative approximation interpretation specified in the main report,
not by quantification over arbitrary sets of real numbers inside iPRA.
